% ----------------------------------------------------------------------------
%                                   RESUMO
% ----------------------------------------------------------------------------

\begin{resumo}[RESUMO]

    \begin{SingleSpacing}
    
Este trabalho apresenta o desenvolvimento de um sistema \textit{web} responsivo (TEia) voltado ao apoio no manejo comportamental de estudantes diagnosticados com Transtorno do Déficit de Atenção com Hiperatividade (TDAH) e Transtorno Opositor Desafiador (TOD). O objetivo central é promover a comunicação integrada entre o corpo docente, as famílias e a gestão escolar, de modo a alinhar estratégias institucionais de acompanhamento e reduzir inconsistências no processo educativo. A pesquisa parte da constatação de que a ausência de comunicação eficaz e os gargalos no registro de ocorrências comprometem o desenvolvimento acadêmico e social desses estudantes, dificultando a aplicação de intervenções consistentes. Metodologicamente, o estudo classifica-se como uma pesquisa aplicada e qualitativa, tendo utilizado instrumentos como entrevistas semiestruturadas, questionários \textit{online} e observação participante na comunidade escolar do Instituto Federal de Brasília (IFB) para o levantamento preciso de requisitos. Distanciando-se de propostas limitadas à concepção teórica, o projeto contemplou a implementação integral da ferramenta, codificada de ponta a ponta com tecnologias \textit{web} nativas (\textit{Vanilla JS}) e infraestrutura em nuvem (Firebase). O sistema entrega funcionalidades de registro ágil de comportamentos, um assistente imediato de estratégias pedagógicas proativas, relatórios institucionais com filtros dinâmicos e controle de acesso rigoroso baseado em perfis. Os resultados consolidados apontam para a otimização da supervisão escolar, a redução da sobrecarga docente em momentos de crise e o fortalecimento do acompanhamento familiar por meio de uma linha do tempo móvel, reafirmando o papel da tecnologia como ponte fundamental para a educação inclusiva.

        \item\textbf{Palavras-chave}: Sistema web. Manejo comportamental. TDAH. TOD. Educação inclusiva.
    
    \end{SingleSpacing}
    
\end{resumo}
